\documentclass[11pt,a4paper]{article}
\usepackage[utf8]{inputenc}
\usepackage[T1]{fontenc}
\usepackage{amsfonts}
\usepackage[french]{babel}
\frenchbsetup{StandardLists=true}
\usepackage{enumitem}
\usepackage{amssymb}
\usepackage{amsmath}
\usepackage[left=2cm,right=2cm,top=2cm,bottom=2cm]{geometry}
\usepackage{multirow}
\usepackage[dvipsnames]{xcolor}
\usepackage{graphicx}
\usepackage{setspace}
\singlespacing
\setlength{\parindent}{0pt}
\usepackage{multicol}
\setlength{\columnseprule}{1pt}
\usepackage{fancyhdr}
\pagestyle{fancy}
\renewcommand\headrulewidth{1pt}
\fancyhead[L]{Lycée Denis Diderot }
\fancyhead[R]{Classe de Terminale}
\setcounter{page}{1}\fancyfoot[C]{page \thepage}
\fancyfoot[R]{Année 2025 2026}
\fancyfoot[L]{Alexis RUIZ - CORRIGÉ}
\usepackage{multirow,tabularx,booktabs}
\usepackage{adjustbox}
\usepackage{tikz}
\newcommand{\cadreblanc}[1]{%
\medskip\framebox[\linewidth][c]{%
\rule{0mm}{#1mm}}\par
\medskip}
\renewcommand{\arraystretch}{2}
\newcommand*\acocher{\quitvmode{\fboxsep0pt \fboxrule0.8pt \fbox{\vrule height1.8ex depth.1ex width0pt \vrule height0pt depth0pt width1.9ex }}}
\usepackage{array}

% Commandes optimisées pour les oscilloscopes du QCM
\newcommand{\oscibase}[1]{%
\begin{tikzpicture}[scale=0.18]
\draw[step=1 cm,gray!60,thin] (-5,-4) grid (5,4);
\draw[thick,-] (0,-4) -- (0,4);
\draw[thick,-] (-5,0) -- (5,0);
#1
\draw[black,very thick,rounded corners=5pt] (-5.1,-4.1) rectangle (5.1,4.1);
\end{tikzpicture}%
}

\newcommand{\oscidoublealternance}{%
\draw[black!70!black,thick,domain=-5:5,samples=50,smooth] 
    plot (\x,{abs(2.5*sin(\x*72))});%
}

\newcommand{\oscimonoalternance}{%
\draw[black!70!black,thick,domain=-5:5,samples=50,smooth] 
    plot (\x,{max(0, 2.5*sin(\x*72))});%
}

\newcommand{\oscicondensateur}{%
\draw[black!70!black,thick,domain=-5:5,samples=80,smooth] 
    plot (\x,{2.3 + 0.2*sin(\x*144)});%
}

\begin{document}
\center \LARGE{Évaluation - Comment redresser une tension alternative ?  }\\[0,3cm]
\normalsize
\hrule

\vspace{2mm}

\flushleft \textbf{NOM :\\[0.2cm]
Prénom:}

\vspace{2mm}

\hrule

\vspace{5mm}

Exercice 1 :
Les oscillogrammes ci-dessous ont été obtenus avec le même balayage (5 ms/div) et la même sensibilité (2 V/div).

\def\ech{0.23}

\begin{center}
\begin{tabular}{|c|c|c|c|c|}
\hline
oscillogramme 1 & oscillogramme 2 & oscillogramme 3 & oscillogramme 4 & oscillogramme 5 \\


 \begin{tikzpicture}[scale=\ech]
\draw[step=1 cm,gray!60,thin] (-5,-4) grid (5,4);
\draw[thick,-] (0,-4) -- (0,4) ;
\draw[thick,-] (-5,0) -- (5,0) ;
\foreach \x in {-5,-4,...,5} {
    \draw (\x,0.1) -- (\x,-0.1);
}
\foreach \y in {-4,-3,...,4} {
    \draw (0.1,\y) -- (-0.1,\y);
}
\draw[thick,-] (-5,2) -- (5,2) ;
\draw[black,very thick,rounded corners=5pt] (-5.1,-4.1) rectangle (5.1,4.1);
\node[gray,font=\footnotesize] at (-4,-5) {2V/div};
\node[gray,font=\footnotesize] at (4,-5) {5ms/div};
\end{tikzpicture}

 &

\begin{tikzpicture}[scale=\ech]
\pgfmathsetmacro{\ampl}{3}
\draw[step=1 cm,gray!60,thin] (-5,-4) grid (5,4);
\draw[thick,-] (0,-4) -- (0,4) ;
\draw[thick,-] (-5,0) -- (5,0) ;
\foreach \x in {-5,-4,...,5} {
    \draw (\x,0.1) -- (\x,-0.1);
}
\foreach \y in {-4,-3,...,4} {
    \draw (0.1,\y) -- (-0.1,\y);
}
\draw[black!70!black,thick] 
    (-5,-\ampl) -- (-1,\ampl)
    (-1,-\ampl) -- (3,\ampl)
    (3,-\ampl) -- (5,-\ampl+2*\ampl*0.5);
\draw[black,very thick,rounded corners=5pt] (-5.1,-4.1) rectangle (5.1,4.1);
\node[gray,font=\footnotesize] at (-4,-5) {2V/div};
\node[gray,font=\footnotesize] at (4,-5) {5ms/div};
\end{tikzpicture}


 &

 
 \begin{tikzpicture}[scale=\ech]
\draw[step=1 cm,gray!60,thin] (-5,-4) grid (5,4);
\draw[thick,-] (0,-4) -- (0,4) ;
\draw[thick,-] (-5,0) -- (5,0) ;
\foreach \x in {-5,-4,...,5} {
    \draw (\x,0.1) -- (\x,-0.1);
}
\foreach \y in {-4,-3,...,4} {
    \draw (0.1,\y) -- (-0.1,\y);
}
\draw[black!70!black,thick,domain=-5:5,samples=50,smooth] 
    plot (\x,{2.5*sin(\x*72)});
\draw[black,very thick,rounded corners=5pt] (-5.1,-4.1) rectangle (5.1,4.1);
\node[gray,font=\footnotesize] at (-4,-5) {2V/div};
\node[gray,font=\footnotesize] at (4,-5) {5ms/div};
\end{tikzpicture}

& 

\begin{tikzpicture}[scale=\ech]
\draw[step=1 cm,gray!60,thin] (-5,-4) grid (5,4);
\draw[thick,-] (0,-4) -- (0,4) ;
\draw[thick,-] (-5,0) -- (5,0) ;
\foreach \x in {-5,-4,...,5} {
    \draw (\x,0.1) -- (\x,-0.1);
}
\foreach \y in {-4,-3,...,4} {
    \draw (0.1,\y) -- (-0.1,\y);
}
\draw[black!70!black,thick] 
    (-5,-3) -- (-2.5,3)
    -- (0,-3)
    -- (2.5,3)
    -- (5,-3);
\draw[black,very thick,rounded corners=5pt] (-5.1,-4.1) rectangle (5.1,4.1);
\node[gray,font=\footnotesize] at (-4,-5) {2V/div};
\node[gray,font=\footnotesize] at (4,-5) {5ms/div};
\end{tikzpicture} 

&

\begin{tikzpicture}[scale=\ech]
\draw[step=1 cm,gray!60,thin] (-5,-4) grid (5,4);
\draw[thick,-] (0,-4) -- (0,4) ;
\draw[thick,-] (-5,0) -- (5,0) ;
\foreach \x in {-5,-4,...,5} {
    \draw (\x,0.1) -- (\x,-0.1);
}
\foreach \y in {-4,-3,...,4} {
    \draw (0.1,\y) -- (-0.1,\y);
}
\draw[black!70!black,thick] 
    (-5,2) -- (-3.75,2)
    -- (-3.75,-2) -- (-2.5,-2)
    -- (-2.5,2) -- (-1.25,2)
    -- (-1.25,-2) -- (0,-2)
    -- (0,2) -- (1.25,2)
    -- (1.25,-2) -- (2.5,-2)
    -- (2.5,2) -- (3.75,2)
    -- (3.75,-2) -- (5,-2);
\draw[black,very thick,rounded corners=5pt] (-5.1,-4.1) rectangle (5.1,4.1);
\node[gray,font=\footnotesize] at (-4,-5) {2V/div};
\node[gray,font=\footnotesize] at (4,-5) {5ms/div};
\end{tikzpicture} \\
\hline
\end{tabular}

\vspace{0.3cm}

\end{center}



\begin{enumerate}
    \item Quelles sont les tensions qui ont la même période ? 
    
    \vspace{3mm}

    \dotfill 
    \item Déterminer alors cette période, T, en secondes, et sa fréquence correspondante, f, en Hertz.
    \begin{multicols}{2}
    \cadreblanc{10}
    \cadreblanc{10}
    \end{multicols}

\item	Quels sont les oscillogrammes qui ont la même tension maximale ?\\ Calculer, en Volt, la valeur de $\mathrm{U_{max}}$.

    \begin{multicols}{2}
    \cadreblanc{10}
    \cadreblanc{10}
    \end{multicols}

   \item 	Pour l'oscillogramme numéro 3, calculer, en V, la tension efficace $\mathrm{U_{eff}}$.

   \cadreblanc{10}

\begin{multicols}{2}
    \item Si pour l'oscillogramme numéro 3 on avait réglé la sensibilité à 1 V/div, quel aurait été son aspect ?
\begin{center}
\begin{tikzpicture}[scale=0.45]
\draw[step=1 cm,gray!60,thin] (-5,-4) grid (5,4);
\draw[thick,-] (0,-4) -- (0,4) ;
\draw[thick,-] (-5,0) -- (5,0) ;
\foreach \x in {-5,-4,...,5} {
    \draw (\x,0.1) -- (\x,-0.1);
}
\foreach \y in {-4,-3,...,4} {
    \draw (0.1,\y) -- (-0.1,\y);
}

\draw[black,very thick,rounded corners=5pt] (-5.1,-4.1) rectangle (5.1,4.1);
\node[gray,font=\footnotesize] at (-4,-5) {1V/div};
\node[gray,font=\footnotesize] at (4,-5) {5ms/div};
\end{tikzpicture}
\end{center}

\columnbreak

\item Si pour l'oscillogramme n°3 on avait réglé le balayage à 10 ms/div (en remettant la sensibilité à 2 V/div), quel aurait été son aspect ?

\begin{center}
\framebox[6cm][c]{\rule{0mm}{40mm}}
\end{center}


    
\end{multicols}
    
\end{enumerate}
\vfill 

\newpage

Exercice 2 : QCM (cocher la ou les bonnes réponses) 

\begin{center}
\begin{tabularx}{\textwidth}{|>{\centering\arraybackslash}m{0.24\textwidth}|>{\centering\arraybackslash}m{0.22\textwidth}|>{\centering\arraybackslash}m{0.22\textwidth}|>{\centering\arraybackslash}m{0.22\textwidth}|}
\hline
\textbf{Le redressement à l'aide d'un pont de diodes permet:} & $\square$ de fournir une tension continue à partir d'une tension alternative. & $\square$ d'obtenir un interrupteur dans le circuit. & $\square$ de fournir une tension alternative à partir d'une tension continue. \\
\hline
\textbf{En observant la diode, le courant circule de:} & $\square$ la droite vers la gauche. & $\square$ la gauche vers la droite. & $\square$ dans les deux sens. \\
\hline
\textbf{Une diode est un composant:} & $\square$ polarisé & $\square$ linéaire & $\square$ possédant trois bornes \\
\hline
\textbf{Le chargeur délivre} & $\square$ du courant continu & $\square$ du courant alternatif & $\square$ de l'énergie chimique \\
\hline
\textbf{Il est possible de recharger:} & $\square$ uniquement des accumulateurs & $\square$ les accumulateurs et les piles & $\square$ uniquement les piles \\
\hline
\textbf{Un accumulateur délivre un courant:} & $\square$ alternatif & $\square$ continu & $\square$ redressé \\
\hline
\textbf{Une diode} & $\square$ autorise le passage du courant dans les deux sens & $\square$ autorise le passage du courant dans un seul sens & $\square$ n'autorise pas le passage du courant \\
\hline
\textbf{Une tension redressée est une tension} & $\square$ qui garde la même valeur & $\square$ qui garde le même signe & $\square$ toujours positive \\
\hline
\textbf{Lorsque l'on réalise le montage correspondant au schéma ci-dessus et que l'on ferme l'interrupteur:} & $\square$ Le courant ne circule pas dans le circuit & $\square$ La lampe ne s'allume pas dans le circuit & $\square$ La lampe s'allume dans le circuit \\
\hline
\textbf{Une tension redressée par un pont de diode correspond:} & 
\adjustbox{valign=m}{\begin{tabular}{c}
\oscibase{\oscidoublealternance}\\[0.3cm]
$\square$
\end{tabular}} & 
\adjustbox{valign=m}{\begin{tabular}{c}
\oscibase{\oscimonoalternance}\\[0.3cm]
$\square$
\end{tabular}} & 
\adjustbox{valign=m}{\begin{tabular}{c}
\oscibase{\oscicondensateur}\\[0.3cm]
$\square$
\end{tabular}} \\
\hline
\textbf{La tension obtenue après un chargeur est:} & 
\adjustbox{valign=m}{\begin{tabular}{c}
\oscibase{\oscidoublealternance}\\[0.3cm]
$\square$
\end{tabular}} & 
\adjustbox{valign=m}{\begin{tabular}{c}
\oscibase{\oscimonoalternance}\\[0.3cm]
$\square$
\end{tabular}} & 
\adjustbox{valign=m}{\begin{tabular}{c}
\oscibase{\oscicondensateur}\\[0.3cm]
$\square$
\end{tabular}} \\
\hline
\end{tabularx}
\end{center}


\end{document}